% Self-contained description of the supplied MATLAB reference model.
% Compile twice with pdflatex; see docs/README.md from the repository root.
\documentclass[11pt,a4paper]{article}
\usepackage[T1]{fontenc}
\usepackage{lmodern}
\usepackage[margin=23mm,headheight=15pt,headsep=8mm,footskip=11mm]{geometry}
\usepackage{amsmath,amssymb,bm,mathtools}
\usepackage{booktabs,array}
\usepackage{xcolor,fancyhdr,hyperref}
\definecolor{navy}{HTML}{17324D}
\definecolor{slate}{HTML}{4C6173}
\hypersetup{colorlinks=true,linkcolor=navy,urlcolor=navy,
pdftitle={Analytical Magnetic Acceleration Model: Reference Equations and Implementation Review},
pdfauthor={SixMagManipulator project},
pdfsubject={Six-magnet MATLAB reference model and analytical Jacobians}}
\pagestyle{fancy}
\fancyhf{}
\fancyhead[L]{\small\textcolor{slate}{SixMagManipulator}}
\fancyhead[R]{\small\textcolor{slate}{Reference model review}}
\fancyfoot[L]{\footnotesize\textcolor{slate}{13 September 2026 \quad | \quad Revision 1}}
\fancyfoot[R]{\footnotesize\thepage}
\renewcommand{\headrulewidth}{0.3pt}
\setlength{\parindent}{0pt}
\setlength{\parskip}{5pt}
\setlength{\emergencystretch}{2em}
\renewcommand{\arraystretch}{1.16}
\numberwithin{equation}{section}
\newcommand{\R}{\mathbb{R}}
\newcommand{\thv}{\bm{\theta}}
\newcommand{\vct}{\operatorname{vec}}
\newcommand{\code}[1]{\texttt{\detokenize{#1}}}
\newcommand{\source}[1]{\textit{Source:} \code{#1}.}
\newcolumntype{L}[1]{>{\raggedright\arraybackslash}p{#1}}

\begin{document}
\thispagestyle{empty}
\begin{center}
{\small\textcolor{slate}{SIXMAGMANIPULATOR \quad / \quad CONTROL MODEL DOCUMENTATION}}\\[8mm]
{\LARGE\bfseries\textcolor{navy}{Analytical Magnetic\\[2mm] Acceleration Model}}\\[5mm]
{\large Reference equations, analytical Jacobians,\\[1mm]
verification results, and implementation review}\\[6mm]
{\small 13 September 2026 \qquad Revision 1}
\end{center}
\vspace{4mm}
\begin{abstract}
\noindent
This report reconstructs the model implemented by the nine MATLAB functions
and two data files in \code{scr/_control/__refmod}. The active model represents
each permanent magnet by 144 fitted point sources, transforms their fields
into the workspace frame, and obtains magnetic acceleration from the gradient
of a scaled field energy. Position and magnet-angle Jacobians are analytical.
Read-only MATLAB checks support the consistency of the active implementation
at sampled configurations. The report records coordinate conventions, required
parameters, derivative formulas, verification results, and opportunities to
avoid repeated calculations. The inactive dipole branch contains derivative
inconsistencies and is documented separately. Physical field units and the
mapping between model angle zero and the GUI pole orientation remain subject
to confirmation.
\end{abstract}

\section{Scope and review status}
\label{sec:scope}
The purpose is to establish a reusable mathematical description before
optimizing or integrating the model into a hardware control loop. The reviewed
files were not modified. The numerical checks in Section~\ref{sec:verification}
verify internal consistency; they do not constitute experimental validation
of the fitted field or bead dynamics.

The acceleration is denoted by $g_n(r,\thv)$, corresponding to $g(r,\theta)$
in the existing \textit{Homotopy method - Analysis} control discussion.
Its Jacobians retain the notation $G_r$ and $G_\theta$.
\begin{center}
\begin{tabular}{L{0.30\linewidth}L{0.62\linewidth}}
\toprule
Item & Reviewed implementation\\
\midrule
Active model & \code{model_sel = 0}: fitted multiple-source model\\
Inputs & $r\in\R^2$ in metres; $\thv\in\R^6$ in radians\\
Outputs & $g_n\in\R^2$, $G_r\in\R^{2\times2}$,
          $G_\theta\in\R^{2\times6}$\\
Geometry & Six identical local source models, rotated and translated\\
Data & 144 source positions and 144 signed coefficients per magnet\\
Checks & MATLAB R2021b; supplied functions and data\\
Open conventions & Unscaled field units and model-to-GUI angle zero\\
\bottomrule
\end{tabular}
\end{center}

\clearpage
\section{Notation, geometry, and coordinate transforms}
\label{sec:geometry}
The bead position and permanent-magnet configuration are
\begin{equation}
r=\begin{bmatrix}r_1\\r_2\end{bmatrix}\in\R^2,\qquad
\thv=\begin{bmatrix}\theta_1&\cdots&\theta_6\end{bmatrix}^{T}\in\R^6.
\label{eq:inputs}
\end{equation}
Subscript $k$ identifies a physical magnet; subscript $j$ identifies an
effective source inside one magnet's local model. All trigonometric
calculations use radians.

For $n=6$, fixed placement angles and global magnet centres are
\begin{equation}
\phi_k=\frac{2\pi(k-1)}{n},\qquad
\rho_k=(a+b)\begin{bmatrix}\cos\phi_k\\\sin\phi_k\end{bmatrix},
\qquad k=1,\ldots,n.
\label{eq:centres}
\end{equation}
The supplied geometry has $a=34.3\,\mathrm{mm}$ and $b=9.75\,\mathrm{mm}$, hence
$\|\rho_k\|=44.05\,\mathrm{mm}$. Magnet 1 lies on positive global $r_1$;
positive global $r_2$ lies between magnets 2 and 3.

Define the orientation, relative global position, and local position by
\begin{equation}
\alpha_k=\phi_k+\theta_k,\qquad p_k=r-\rho_k,\qquad
w_k=R_kp_k,\qquad R_k=R(\alpha_k),
\label{eq:local-position}
\end{equation}
where the actual matrix in \code{R.m} is
\begin{equation}
R(\alpha)=
\begin{bmatrix}\cos\alpha&\sin\alpha\\-\sin\alpha&\cos\alpha\end{bmatrix}.
\label{eq:R}
\end{equation}
Here $R_k$ transforms global coordinates into the local frame; $R_k^T$
transforms local vectors into the global frame. Increasing $\theta_k$ rotates
the physical local axes counterclockwise. The comment in \code{R.m} labels
the transform in the opposite direction; its matrix and usage determine
the convention in this report.

The angle derivative is
\begin{equation}
R_{\theta,k}:=\frac{\partial R_k}{\partial\theta_k}
=\begin{bmatrix}-\sin\alpha_k&\cos\alpha_k\\
-\cos\alpha_k&-\sin\alpha_k\end{bmatrix},
\qquad \frac{\partial w_k}{\partial\theta_k}=R_{\theta,k}p_k.
\label{eq:R-derivative}
\end{equation}
$R_{\theta,k}$ is an angle derivative, not a time derivative.
The planar projection and embedding matrices are
\begin{equation}
E=\begin{bmatrix}1&0&0\\0&1&0\end{bmatrix},\qquad
E^T:\R^2\longrightarrow\R^3.
\label{eq:projection}
\end{equation}
Embedding a point as $E^Tw=[w_1,w_2,0]^T$ retains three-dimensional distances
to the sources; it does not discard their heights.
\source{parametersGen.m, R.m, hn_field.m, Hn.m}

\section{Single-magnet field and analytical derivatives}
\label{sec:local-field}
\subsection{Fitted multiple-source field}
Let $\rho_j^c\in\R^3$ and $a_j\in\R$ be the local source positions and
coefficients in \code{par.rho} and \code{par.ak_comp}. Superscript $c$
distinguishes these positions from the global centres $\rho_k$.
For $N=144$, the field at a local point $p\in\R^3$ is
\begin{equation}
\boxed{h_c(p)=\sum_{j=1}^{N}a_j
\frac{p-\rho_j^c}{\|p-\rho_j^c\|^3}.}
\label{eq:hc}
\end{equation}
These are effective fitted sources, not physical free magnetic monopoles.
The code's name \textit{complex model} means this multiple-source
representation. All stored positions and coefficients are real-valued doubles.
The planar field used subsequently is
\begin{equation}
\bar h_c(w):=E h_c(E^Tw),\qquad w\in\R^2.
\label{eq:hc-planar}
\end{equation}
Although the header of \code{h_c.m} mentions a $2\times1$ input, the active
implementation is called with a $3\times1$ input and returns a $3\times1$
field. Its optional second output is $\|h_c(p)\|$.

\subsection{First spatial derivative}
For each source, define
\begin{equation}
q_j=E^Tw-\rho_j^c,\qquad u_j=Eq_j,\qquad d_j=\|q_j\|.
\label{eq:source-offset}
\end{equation}
For $d_j>0$, differentiation gives
\begin{equation}
\boxed{H_c(w):=\frac{\partial\bar h_c}{\partial w}
=\sum_{j=1}^{N}a_j
\left(\frac{I_2}{d_j^3}-\frac{3u_ju_j^T}{d_j^5}\right).}
\label{eq:Hc}
\end{equation}
This is the upper-left planar block of the three-dimensional Jacobian
assembled in \code{Hc.m}. Each summand is symmetric, so $H_c=H_c^T$.

\subsection{Second spatial derivative}
The derivative tensor of $H_c$ has components
\begin{align}
\frac{\partial(H_c)_{ab}}{\partial w_\ell}
=\sum_{j=1}^{N}a_j\bigg[
&-\frac{3}{d_j^5}
\left(\delta_{ab}u_{j\ell}+\delta_{a\ell}u_{jb}
+\delta_{b\ell}u_{ja}\right)\notag\\
&+\frac{15u_{ja}u_{jb}u_{j\ell}}{d_j^7}\bigg],
\qquad a,b,\ell\in\{1,2\}.
\label{eq:Tc-components}
\end{align}
$\delta_{ab}$ is the Kronecker delta. The function \code{dvecHcdw.m} returns
\begin{equation}
D_c(w):=\frac{\partial\,\vct H_c(w)}{\partial w}\in\R^{4\times2},
\quad \vct H_c=
\begin{bmatrix}
(H_c)_{11}&(H_c)_{21}&(H_c)_{12}&(H_c)_{22}
\end{bmatrix}^{T}.
\label{eq:Dc}
\end{equation}
Vectorization follows MATLAB's column-wise convention.
Define the directional derivative
\begin{equation}
\mathcal T_c(w)[v]:=
\sum_{\ell=1}^{2}\frac{\partial H_c(w)}{\partial w_\ell}v_\ell,
\qquad \vct\bigl(\mathcal T_c(w)[v]\bigr)=D_c(w)v.
\label{eq:directional-T}
\end{equation}
An equivalent expression using two-dimensional vectors is
\begin{align}
\mathcal T_c(w)[v]=\sum_{j=1}^{N}a_j\bigg[
&-\frac{3}{d_j^5}\left((u_j^Tv)I_2+vu_j^T+u_jv^T\right)\notag\\
&+\frac{15(u_j^Tv)}{d_j^7}u_ju_j^T\bigg].
\label{eq:Tc-directional}
\end{align}
This exposes the analytical operations without constructing a
three-dimensional tensor. All formulas assume the evaluation point does not
coincide with an effective source.
\source{h_c.m, Hc.m, dvecHcdw.m}

\section{Combined field, scaling, and magnetic acceleration}
\label{sec:acceleration}
The force-per-mass coefficient is
\begin{equation}
k_g=\frac{\chi}{2\mu_0\rho_{\mathrm{Fe}}(1+\chi/3)},
\qquad s=\sqrt{2k_g}.
\label{eq:kg}
\end{equation}
The supplied values are $\chi=1000$, $\mu_0=4\pi\times10^{-7}$ in SI units,
and $\rho_{\mathrm{Fe}}=7800\,\mathrm{kg/m^3}$, giving
$k_g\approx152.5759$. The force coefficient satisfies $k_m=m_pk_g$, where
$m_p$ is particle mass; $m_p$ therefore cancels from $k_g$ in this model.

Define the unscaled planar combined field as
\begin{equation}
h_{\mathrm{raw}}(r,\thv)=\sum_{k=1}^{n}R_k^T\bar h_c(w_k).
\label{eq:raw-field}
\end{equation}
The output of \code{hn_field.m} is the scaled field
\begin{equation}
\boxed{h_n(r,\thv)=s\,h_{\mathrm{raw}}(r,\thv).}
\label{eq:hn}
\end{equation}
Its position Jacobian, returned by \code{Hn.m}, is
\begin{equation}
\boxed{H_n(r,\thv)=\frac{\partial h_n}{\partial r}
=s\sum_{k=1}^{n}R_k^TH_c(w_k)R_k.}
\label{eq:Hn}
\end{equation}
Equation~\eqref{eq:Hc} implies $H_n=H_n^T$ for the active model.
The acceleration in \code{gn_field.m} is
\begin{equation}
\boxed{g_n(r,\thv)=H_n(r,\thv)h_n(r,\thv).}
\label{eq:gn}
\end{equation}
For a general vector field, the gradient of half its squared norm is
$H_n^Th_n$. Symmetry gives
\begin{equation}
g_n=\nabla_r\left(\frac12\|h_n\|^2\right)
=k_g\nabla_r\|h_{\mathrm{raw}}\|^2.
\label{eq:energy-gradient}
\end{equation}
Thus the implementation uses the gradient of the \emph{squared} field norm.
There is no division by $\|h_n\|$.

Only the planar field enters Equation~\eqref{eq:energy-gradient}. A full
three-dimensional field-energy gradient would also contain the third
component's contribution. Neglecting that component is part of the
transverse-plane model and should be retained deliberately in a translation.

These routines provide magnetic acceleration only. In a dynamics model
\begin{equation}
\dot r=v,\qquad \dot v=-\sigma_vv+g_n(r,\thv),
\label{eq:translation-context}
\end{equation}
drag is evaluated separately; \code{gn_field.m} does not apply
\code{par.sigma_v}.

\source{parametersGen.m, hn_field.m, Hn.m, gn_field.m}

\section{Analytical acceleration Jacobians}
\label{sec:jacobians}
\subsection{Chain rule in matrix form}
Define
\begin{equation}
G_r:=\frac{\partial g_n}{\partial r}\in\R^{2\times2},\qquad
G_\theta:=\frac{\partial g_n}{\partial\thv}\in\R^{2\times6}.
\label{eq:jac-definitions}
\end{equation}
Let $e_\ell$ be column $\ell$ of $I_2$ and let
$v_{\theta,k}=R_{\theta,k}p_k$. Only magnet $k$ depends on $\theta_k$, so
\begin{equation}
\frac{\partial h_n}{\partial\theta_k}
=s\left[R_{\theta,k}^T\bar h_c(w_k)
+R_k^TH_c(w_k)v_{\theta,k}\right].
\label{eq:dh-angle}
\end{equation}
The first term rotates the field vector; the second accounts for the changing
local evaluation point. The derivative of the combined field Jacobian is
\begin{align}
\frac{\partial H_n}{\partial\theta_k}
=s\big[
&R_{\theta,k}^TH_c(w_k)R_k
+R_k^TH_c(w_k)R_{\theta,k}\notag\\
&+R_k^T\mathcal T_c(w_k)[v_{\theta,k}]R_k\big].
\label{eq:dH-angle}
\end{align}
For position derivatives, the orientations are fixed:
\begin{equation}
\frac{\partial H_n}{\partial r_\ell}
=s\sum_{k=1}^{n}R_k^T\mathcal T_c(w_k)[R_ke_\ell]R_k,
\qquad \ell=1,2.
\label{eq:dH-position}
\end{equation}
The product rule applied to Equation~\eqref{eq:gn} yields
\begin{equation}
\boxed{(G_\theta)_{:,k}
=\frac{\partial H_n}{\partial\theta_k}h_n
+H_n\frac{\partial h_n}{\partial\theta_k}.}
\label{eq:Gtheta-columns}
\end{equation}
Similarly,
\begin{equation}
\boxed{G_r=H_n^2+
\begin{bmatrix}
\dfrac{\partial H_n}{\partial r_1}h_n&
\dfrac{\partial H_n}{\partial r_2}h_n
\end{bmatrix}.}
\label{eq:Gr-columns}
\end{equation}
All derivatives in these expressions are analytical.

\subsection{Correspondence with the vectorized implementation}
Introduce
\begin{equation}
B_\theta=\frac{\partial h_n}{\partial\thv},\qquad
D_r=\frac{\partial\,\vct H_n}{\partial r},\qquad
D_\theta=\frac{\partial\,\vct H_n}{\partial\thv}.
\label{eq:vectorized-definitions}
\end{equation}
Their shapes are $2\times6$, $4\times2$, and $4\times6$, respectively.
The final two lines of \code{gan_linear.m} are
\begin{align}
\boxed{G_r=H_n^2+(I_2\otimes h_n^T)D_r,}
\label{eq:Gr-vectorized}\\
\boxed{G_\theta=H_nB_\theta+(I_2\otimes h_n^T)D_\theta.}
\label{eq:Gtheta-vectorized}
\end{align}
For a $2\times2$ matrix $A$, $(I_2\otimes h^T)\vct A=A^Th$. Derivatives of
$H_n$ are symmetric for the active model, making these contractions identical
to Equations~\eqref{eq:Gtheta-columns} and \eqref{eq:Gr-columns}.

The rotation-vector derivatives use the same column-wise convention:
\begin{align}
\frac{\partial\,\vct R_k}{\partial\theta_k}
&=\begin{bmatrix}
-\sin\alpha_k&-\cos\alpha_k&\cos\alpha_k&-\sin\alpha_k
\end{bmatrix}^{T},
\label{eq:vec-R-derivative}\\
\frac{\partial\,\vct R_k^T}{\partial\theta_k}
&=\begin{bmatrix}
-\sin\alpha_k&\cos\alpha_k&-\cos\alpha_k&-\sin\alpha_k
\end{bmatrix}^{T}.
\label{eq:vec-RT-derivative}
\end{align}
For later control development,
\begin{equation}
\mathrm{d}g_n=G_r\,\mathrm{d}r+G_\theta\,\mathrm{d}\thv,\qquad
\dot g_n=G_r\dot r+G_\theta\dot\thv.
\label{eq:control-chain-rule}
\end{equation}
These identities describe sensitivity of the acceleration map; they do not
specify a control law.
\source{gan_linear.m, dvecHcdw.m}

\section{Parameter dependencies and units}
\label{sec:parameters}
With the complex model selected, the numerical kernel directly needs:
\begin{center}
\begin{tabular}{L{0.27\linewidth}L{0.16\linewidth}L{0.47\linewidth}}
\toprule
Parameter & Shape & Purpose\\
\midrule
\code{n} & scalar, 6 & Physical magnet count\\
\code{phi_rad} & $6\times1$ & Fixed placement angles\\
\code{rho_k_rad} & $6\times2$ & Global magnet centres stored by row\\
\code{kg} & scalar & Field-to-acceleration coefficient\\
\code{rho} & $3\times144$ & Local effective source positions\\
\code{ak_comp} & $144\times1$ & Signed effective source coefficients\\
\bottomrule
\end{tabular}
\end{center}
\code{model_sel} is also consulted by the current functions. It can become a
configuration choice if separate model implementations are introduced.
The optional dipole branch additionally uses \code{m} and \code{ak_dipole}.

The constants $a,b,\mu_0,\chi,\rho_{\mathrm{Fe}}$ are needed indirectly during
initialization of the centres and $k_g$. Controller gains, trajectories,
actuator dynamics, \code{delta_r}, \code{delta_th}, and drag parameters are not
used by the active field/Jacobian calculation. Particle radius and mass do
not appear after $k_g$ is formed under the supplied material model.

The following dimensions apply if the unscaled field is magnetic flux density
in tesla, as suggested by Equation~\eqref{eq:kg}:
\begin{center}
\begin{tabular}{ll}
\toprule
Quantity & Units\\
\midrule
$r,\rho_k,\rho_j^c$ & $\mathrm{m}$\\
$\theta_k,\phi_k$ & $\mathrm{rad}$\\
$h_c,\bar h_c,h_{\mathrm{raw}}$ & $\mathrm{T}$\\
$a_j$ & $\mathrm{T\,m^2}$\\
$k_g$ & $\mathrm{m^2\,s^{-2}\,T^{-2}}$\\
$h_n,\ H_n$ & $\mathrm{m\,s^{-1}},\ \mathrm{s^{-1}}$\\
$g_n,\ G_r$ & $\mathrm{m\,s^{-2}},\ \mathrm{s^{-2}}$\\
$G_\theta$ & $\mathrm{m\,s^{-2}\,rad^{-1}}$\\
\bottomrule
\end{tabular}
\end{center}
The scaled field $h_n$ is not dimensionless. Confirmation of the unscaled
field units is required before accepting the physical interpretation.

\section{Numerical consistency checks}
\label{sec:verification}
Read-only checks used the supplied functions in MATLAB R2021b with
\code{par = parametersGen()} and \code{par.model_sel = 0}. Central differences
served only as a verification tool; they are not part of the model evaluation.

With \code{rng(17)}, twelve configurations were generated in this order:
\begin{equation}
r^{(i)}=0.025\sqrt{\eta_i}
\begin{bmatrix}\cos(2\pi i/12)\\\sin(2\pi i/12)\end{bmatrix},\qquad
\thv^{(i)}=2.5(2\xi_i-\mathbf 1).
\label{eq:validation-samples}
\end{equation}
For each $i=1,\ldots,12$, $\eta_i$ is one call to \code{rand} and $\xi_i$ is
the following call to \code{rand(6,1)}. Positions lie within a
$25\,\mathrm{mm}$ radius; magnet angles lie in $[-2.5,2.5]$ radians.

The finite-difference columns were
\begin{align}
(G_r^{\mathrm{FD}})_{:,\ell}
&=\frac{g_n(r+\epsilon_re_\ell,\thv)
-g_n(r-\epsilon_re_\ell,\thv)}{2\epsilon_r},
\label{eq:fd-r}\\
(G_\theta^{\mathrm{FD}})_{:,k}
&=\frac{g_n(r,\thv+\epsilon_\theta e_k^{(6)})
-g_n(r,\thv-\epsilon_\theta e_k^{(6)})}{2\epsilon_\theta}.
\label{eq:fd-theta}
\end{align}
Here $e_k^{(6)}$ is column $k$ of $I_6$, $\epsilon_r=10^{-7}\,\mathrm{m}$,
and $\epsilon_\theta=10^{-6}\,\mathrm{rad}$.
Relative discrepancy was measured by
\begin{equation}
\varepsilon_A=
\frac{\|A_{\mathrm{analytic}}-A_{\mathrm{FD}}\|_F}
{\max(\|A_{\mathrm{FD}}\|_F,\varepsilon_{\mathrm{mach}})}.
\label{eq:relative-error}
\end{equation}
\begin{center}
\begin{tabular}{L{0.15\linewidth}L{0.47\linewidth}r}
\toprule
Quantity & Check domain & Relative discrepancy\\
\midrule
$G_r$ & Maximum over 12 configurations & $1.63\times10^{-10}$\\
$G_\theta$ & Maximum over 12 configurations & $7.57\times10^{-10}$\\
$H_c$ & $w=[-0.04,\ 0.006]^T\,\mathrm{m}$ & $4.08\times10^{-11}$\\
$D_c$ & Same local point & $5.34\times10^{-11}$\\
\bottomrule
\end{tabular}
\end{center}
The local checks differentiated $\bar h_c$ and $\vct H_c$ using the same
spatial step. They are single-point checks, not twelve distinct local points.
These results support internal derivative consistency at the tested points;
they do not establish a uniform error bound, physical field accuracy, or
hardware tracking performance.

\section{Implementation findings and opportunities}
\label{sec:implementation}
\subsection{Repeated work in the active model}
One call to \code{gan_linear} evaluates \code{Hn} and \code{hn_field}, then
runs another per-magnet loop containing \code{Hc}, \code{h_c}, and
\code{dvecHcdw}. Consequently:
\begin{itemize}
\item \code{Hc} and \code{h_c} are each evaluated 12 times, and
      \code{dvecHcdw} is evaluated six times.
\item Source offsets, norms, rotations, and inverse-distance powers are
      recalculated in multiple helpers.
\item \code{Hc} builds a $144\times144$ diagonal matrix for a weighted
      outer-product sum; direct weighted accumulation avoids that matrix.
\item The angle derivatives build identity matrices, selector vectors,
      Kronecker products, and intermediate arrays for small fixed outputs.
\item \code{h_c} calculates the norm output even when only the field is used.
\end{itemize}
A shared analytical evaluator returning $g_n,G_r,G_\theta$ can reuse these
quantities. No speedup has yet been measured. The approximately
$7\,\mathrm{ms}$ Jacobian time was reported by the user and has not been
re-benchmarked or separated into individual costs in this review.

Source positions and coefficients can be loaded once, and constants such as
$\sqrt{2k_g}$ can be formed during initialization. Fixed dimensions are
available for the supplied six-magnet, 144-source configuration.

\subsection{Inactive dipole branch}
The dipole field in \code{h_c.m} is
\begin{equation}
h_{\mathrm{dip}}(p)=a_d\left[
\frac{3(m^Tp)p}{\|p\|^5}-\frac{m}{\|p\|^3}\right],
\qquad m=\begin{bmatrix}1&0&0\end{bmatrix}^{T},
\label{eq:dipole-field}
\end{equation}
with $a_d\approx4.8169\times10^{-7}$ from \code{dipole_model_dat.mat}.
For the planar dipole, $\bar m=Em$, its first derivative should be
\begin{equation}
H_{\mathrm{dip}}(w)=3a_d\left[
\frac{w\bar m^T+\bar m w^T+(\bar m^Tw)I_2}{\|w\|^5}
-\frac{5(\bar m^Tw)ww^T}{\|w\|^7}\right].
\label{eq:dipole-correct-derivative}
\end{equation}
The existing \code{Hc.m} dipole branch is not this derivative. At
$w=[-0.04,\ 0.006]^T\,\mathrm{m}$, its relative discrepancy from a central
difference of \code{h_c} was approximately $0.385$, and its Jacobian was not
symmetric. The \code{model_sel = 1} branch of \code{dvecHcdw.m} also uses the
complex-model source data.

Equation~\eqref{eq:dipole-correct-derivative} records the mathematical
comparison; no fix has been applied. The default \code{model_sel = 0} does
not select these inconsistent branches.

\section{Decisions before control-loop integration}
\label{sec:decisions}
\begin{enumerate}
\item \textbf{Model scope.} Confirm whether production work should initially
      support only the verified complex-model branch.
\item \textbf{Angle zero and pole direction.} At $\theta_k=0$, local positive
      $x$ points radially outward. The fitted source first moment is mainly
      along local positive $x$. The GUI depicts North inward at zero;
      confirm the intended physical mapping and any angular offset.
\item \textbf{Physical field units.} Confirm that the unscaled \code{h_c}
      output is magnetic flux density $B$ in tesla. The coefficient $k_g$
      supports this interpretation rather than field strength in amperes
      per metre.
\item \textbf{Timing budget.} Confirm whether the reported $7\,\mathrm{ms}$
      covers only \code{gan_linear} with a preinitialized parameter structure,
      and specify how many field/Jacobian evaluations one control update needs.
\end{enumerate}
The model applies the same fitted sources to all six magnets.
Per-magnet differences, if required by calibration, need an explicit extension
of the data representation.

The GUI reports positions in millimetres and motor angles in degrees.
The model boundary must convert
\begin{equation}
r_{\mathrm{m}}=10^{-3}r_{\mathrm{mm}},\qquad
\thv_{\mathrm{rad}}=\frac{\pi}{180}\thv_{\mathrm{deg}},
\label{eq:unit-boundary}
\end{equation}
in addition to the confirmed coordinate and angle-zero transforms.
Sensitivities with respect to displayed units follow the chain rule:
\begin{equation}
\frac{\partial g_n}{\partial r_{\mathrm{mm}}}=10^{-3}G_r,\qquad
\frac{\partial g_n}{\partial\thv_{\mathrm{deg}}}
=\frac{\pi}{180}G_\theta.
\label{eq:jacobian-unit-boundary}
\end{equation}
Any reversal of the displayed workspace axes must be handled consistently
when converting to the fixed model frame.

Servo bias correction also belongs at the interface: the controller uses
magnet angles relative to the configured zero, and the corresponding bias is
added once when converting a commanded angle into a raw servo goal. The
unresolved pole-orientation offset must not be silently confused with encoder
bias.

\appendix
\section{Reference data and source traceability}
\label{sec:provenance}
The supplied MAT files contain:
\begin{center}
\begin{tabular}{L{0.47\linewidth}L{0.44\linewidth}}
\toprule
Quantity & Value\\
\midrule
Complex-model source array & $3\times144$ real doubles\\
Complex-model coefficient array & $144\times1$ real doubles\\
Source $x$ and $y$ extents & Approximately $\pm8.8271\,\mathrm{mm}$\\
Source $z$ extent & $\pm9.5\,\mathrm{mm}$\\
Distinct source $z$ levels & Nine; spacing $2.375\,\mathrm{mm}$\\
Coefficient minimum / maximum &
  $-1.0802140\times10^{-6}$ / $1.0818149\times10^{-6}$\\
Coefficient sum & $5.0963\times10^{-10}$\\
\bottomrule
\end{tabular}
\end{center}
The local first moment is
\begin{equation}
\sum_{j=1}^{144}a_j\rho_j^c
\approx\begin{bmatrix}
4.97344\times10^{-7}\\
-2.66832\times10^{-12}\\
-6.73725\times10^{-12}
\end{bmatrix}.
\label{eq:source-first-moment}
\end{equation}
The nonzero coefficient sum is recorded as supplied; this review does not
impose exact antisymmetry or alter fitted values.

\begin{center}
\small
\begin{tabular}{L{0.32\linewidth}L{0.33\linewidth}L{0.25\linewidth}}
\toprule
File & Role & SHA-256 prefix\\
\midrule
\code{parametersGen.m} & Initialization and data loading & \code{af5187bc3eff51f8}\\
\code{R.m} & Coordinate rotation & \code{a163061a87037f02}\\
\code{h_c.m} & Local three-dimensional field & \code{59a62c8f04bcaed4}\\
\code{Hc.m} & Local planar field Jacobian & \code{b2b00644eb4a0dc0}\\
\code{dvecHcdw.m} & Local second spatial derivative & \code{14799653b044de62}\\
\code{hn_field.m} & Scaled combined planar field & \code{8f8199a5b1e02746}\\
\code{Hn.m} & Combined field Jacobian & \code{aa62f606972fe9a2}\\
\code{gn_field.m} & Magnetic acceleration & \code{6ec10d8a7de93aa1}\\
\code{gan_linear.m} & Acceleration Jacobians & \code{9bc1221ad9bd593c}\\
\code{complex_model_dat.mat} & Fitted sources and coefficients & \code{272ce8c57cb58148}\\
\code{dipole_model_dat.mat} & Dipole coefficient & \code{0790869d7af10a17}\\
\bottomrule
\end{tabular}
\end{center}
Fingerprints are the first 16 hexadecimal digits of each file's SHA-256 hash.
They were computed on the report date from \code{scr/_control/__refmod}. The filenames,
equations, and fingerprints identify this reference independently of later
optimization work.

\section{Reuse in program documentation}
Sections~\ref{sec:geometry}--\ref{sec:jacobians} form the mathematical model
description. Section~\ref{sec:parameters} specifies kernel data and units.
Sections~\ref{sec:verification}--\ref{sec:decisions} record the current evidence,
implementation findings, and unresolved decisions; they should be updated when
native implementation or physical calibration is introduced.

The LaTeX source is self-contained: it has no external figures, bibliography
database, or generated equation fragments. Equation labels may be retained
when these sections are incorporated into a larger software description.
\end{document}
